\section*{Chapitre \ref{ch:b2:precambrian-earth} --- La Terre précambrienne}

\begin{solution}{exo:b2:precambrian-earth:1}
(a) \omterm{def:b2:precambrian-earth:eons}{Hadéen} $4567 - 4031 = \qty{536}{Ma}$~; \omterm{def:b2:precambrian-earth:eons}{Archéen} $4031 - 2500 = \qty{1531}{Ma}$~;
\omterm{def:b2:precambrian-earth:eons}{Protérozoïque} $2500 - 538.8 = \qty{1961}{Ma}$. (b) $(4567 - 538.8)/4567 = 0.882$~: environ sept
huitièmes. (c) $4567 - 4404 = \qty{163}{Ma}$ pour le zircon~; $4567 - 4030 = \qty{537}{Ma}$ pour
les roches les plus anciennes.
\end{solution}

\begin{solution}{exo:b2:precambrian-earth:2}
(a) Les facteurs $e^{\lambda t}$ avec $\lambda = \ln2/T$ et $t = \qty{3}{Ga}$ valent 1.593
(\ce{^{238}U}), 19.18 (\ce{^{235}U}), 1.160 (\ce{^{232}Th}) et 5.292 (\ce{^{40}K})~: les
contributions valent 3.053, 1.573, 2.424 et \qty{4.329}{pW}, soit \qty{11.38}{pW} au total, contre
\qty{4.906}{pW} aujourd'hui. (b) $11.38/4.906 = 2.32$. (c) Le thorium-232 domine aujourd'hui~; le potassium-40 dominait
alors.
\end{solution}

\begin{solution}{exo:b2:precambrian-earth:3}
(a) $1616 + 18.0\times(28 - 32.2) = \qty{1540}{\celsius}$. (b) $32.2 + (1500 - 1616)/18.0 =
\qty{25.8}{\percent}$. (c) Une analyse de roche totale n'est pas la composition d'un liquide~: une
coulée de \omterm{def:b2:precambrian-earth:greenstone}{komatiite} peut contenir des cristaux d'olivine accumulés à partir du liquide, qui ajoutent du MgO. Son
\qty{38}{\percent} dépasse les \qty{32}{\percent} des liquides de \omterm{def:b2:precambrian-earth:greenstone}{komatiite} les plus riches, et
les expériences sur lesquelles repose la droite.
\end{solution}

\begin{solution}{exo:b2:precambrian-earth:4}
Roches les plus anciennes, \qty{4.03}{Ga}~; premiers stromatolithes, \qty{3.43}{Ga}~; fin des dômes de la Pilbara,
\qty{3.1}{Ga}~; \omterm{def:b2:precambrian-earth:goe}{Grande Oxydation}, environ 2.4 à \qty{2.1}{Ga}~; dernières \omterm{def:b2:precambrian-earth:greenstone}{komatiites} courantes,
\qty{1.9}{Ga}~; seconde montée de l'oxygène, environ \qty{0.8}{Ga}~; fin de la glaciation marinoenne,
\qty{0.635}{Ga}~; premiers fossiles édiacariens, environ \qty{0.57}{Ga}.
\end{solution}

\begin{solution}{exo:b2:precambrian-earth:5}
(a) $\tfrac14\times1361\times0.706 = \qty{240.2}{W.m^{-2}}$ absorbés et $\sigma\times288^4 =
\qty{390.1}{W.m^{-2}}$, donc $\varepsilon = 240.2/390.1 = 0.616$. (b) $\tfrac14\times1361\times0.4
= \qty{136.1}{W.m^{-2}}$, $T^4 = 136.1/(0.616\times\num{5.670e-8})$, $T = \qty{249.8}{K}$.
(c) À \qty{275.5}{K}, l'albédo vaut $0.6 - 0.306\times12.5/20 = 0.4088$~: $R = 340.25\times0.5912 =
\qty{201.2}{W.m^{-2}}$ et $E = 0.616\,\sigma\times275.5^4 = \qty{201.2}{W.m^{-2}}$. (d) $R' =
340.25\times0.306/20 = \qty{5.21}{W.m^{-2}.K^{-1}}$, $E' = 4\varepsilon\sigma T^3 =
\qty{2.92}{W.m^{-2}.K^{-1}}$~: $R' > E'$, donc l'équilibre est instable
(\cref{prop:b2:precambrian-earth:ebm}).
\end{solution}

\begin{solution}{exo:b2:precambrian-earth:6}
(a) $1/(1 + 0.4\times2.4/4.567) = 0.826$~; $1/(1 + 0.4\times0.65/4.567) = 0.946$. (b) L'état chaud
exige au moins 0.932 du rayonnement actuel~: impossible il y a \qty{2.4}{Ga}, possible
il y a \qty{0.65}{Ga}. (c) $0.826/0.946 = 0.873$~: \qty{13}{\percent} moins brillant, proche de la
valeur de \qty{10}{\percent} citée, qui vient d'un autre modèle solaire.
\end{solution}

\begin{solution}{exo:b2:precambrian-earth:7}
(a) Sturtien $716.9 - 658.9 = \qty{58.0}{Ma}$~; Marinoen $644.5 - 635.4 = \qty{9.1}{Ma}$~;
interglaciaire $658.9 - 644.5 = \qty{14.4}{Ma}$. (b) $(58.0 + 9.1)/85 = 0.79$~: environ quatre cinquièmes
de la période. (c) Son début n'est encadré qu'entre 649.9 et \qty{639.0}{Ma}~: avec ces
bornes, le Marinoen a duré entre environ 3 et \qty{15}{Ma}.
\end{solution}

\begin{solution}{exo:b2:precambrian-earth:8}
(a) $\tau = \frac{0.05}{0.95}\times\frac{44.01}{28.97}\times\frac{\num{5.1e18}}{\num{0.26e12}} =
0.05263\times1.519\times\num{1.962e7} = \qty{1.57}{Ma}$. (b) $0.05/\num{278e-6} = 180$ fois.
(c) Deux fois plus long, \qty{3.14}{Ma}.
\end{solution}

\begin{solution}{exo:b2:precambrian-earth:9}
(a) Les pierres sont tombées dans une boue molle et en ont courbé les lamines vers le bas~: ce sont des pierres lâchées (\textit{dropstones}), apportées par
la glace flottante. (b) Une glace atteignant la mer à \ang{10} de l'équateur est la trace attendue d'une
boule de neige, et très difficile à obtenir si seules les régions polaires sont gelées. (c) Un \omterm{def:b2:precambrian-earth:snowball}{carbonate de couverture}~:
il enregistre la fin brutale de la glaciation, avec une altération intense et l'inondation des
plateaux par l'eau de fonte.
\end{solution}

\begin{solution}{exo:b2:precambrian-earth:10}
(a) $0.94 > 0.932$, le point de bascule chaud~: l'équilibre chaud, $288\times0.94^{1/4} =
\qty{283.6}{K}$, se place juste au-dessus de \qty{283}{K}~; et $0.94 < 1.227$, donc la boule de neige existe aussi.
(b) $249.8\times0.94^{1/4} = \qty{246.0}{K}$. (c) $\varepsilon = 0.94\times1361\times0.4/(4\sigma
\times263^4) = 511.7/1085.1 = 0.472$. (d) $\sigma\times263^4 = \qty{271.3}{W.m^{-2}}$, et $(0.616
- 0.472)\times271.3 = \qty{39}{W.m^{-2}}$.
\end{solution}

\begin{solution}{exo:b2:precambrian-earth:11}
(a) $S_{\text{m}}/S_{\text{f}} = (T_1/T_2)^4\,(1 - a_{\text{w}})/(1 - a_{\text{i}})$~: l'émissivité
s'annule. $(263/283)^4\times0.706/0.4 = 0.7459\times1.765 = 1.317$. (b) Un contraste
d'albédo plus faible entre la glace et l'eau libre, comme une glace plus sombre et poussiéreuse, ou une rampe de
température plus large. (c) Une boule de neige réfléchit tant de lumière qu'elle reste froide sous tout rayonnement solaire ou
effet de serre auquel une Terre chaude gèlerait~; pour la faire dégeler, le forçage doit dépasser celui du
gel, sur toute la largeur de la boucle.
\end{solution}

\begin{solution}{exo:b2:precambrian-earth:12}
Près de $T^*$, $C\dot\delta = -(E' - R')\delta$, donc $\delta \propto e^{-t/\tau}$ avec $\tau =
C/(E' - R')$. À l'équilibre chaud, $R' = 0$ et $E' = 4\times0.616\times\num{5.670e-8}\times
288^3 = \qty{3.34}{W.m^{-2}.K^{-1}}$; $C = 1025\times3992\times50 = \qty{2.05e8}{J.m^{-2}.K^{-1}}$,
donc $\tau = \qty{6.1e7}{s}$, soit environ \qty{1.9}{an}. Dans ce modèle, l'équilibre chaud garde $R' =
0$ jusqu'à atteindre \qty{283}{K}, si bien que $\tau$ reste voisin de deux ans jusqu'au point de bascule, puis
l'état disparaît. Avec un albédo variant continûment, $R'$ croît à mesure que l'équilibre approche du
repli, où $E' = R'$~: $\tau$ croît sans borne, un ralentissement qui annonce le point de
bascule.
\end{solution}

\begin{solution}{pb:b2:precambrian-earth:1}
\textbf{1.} $\num{0.0203e-6}\times0.9929\times\num{95.13e-6} = \qty{1.92}{pW.kg^{-1}}$
(\ce{^{238}U}); $\num{0.0203e-6}\times0.00711\times\num{568.47e-6} = \qty{0.0820}{pW.kg^{-1}}$
(\ce{^{235}U}); $\num{0.0795e-6}\times\num{26.28e-6} = \qty{2.09}{pW.kg^{-1}}$ (\ce{^{232}Th});
$\num{240e-6}\times\num{1.20e-4}\times\num{28.47e-6} = \qty{0.820}{pW.kg^{-1}}$ (\ce{^{40}K}).
\textbf{2.} \qty{4.91}{pW.kg^{-1}}.
\textbf{3.} $\ln2/T$~: 0.155, 0.985, 0.0495 et \qty{0.555}{Ga^{-1}}.
\textbf{4.} Facteurs 1.52, 14.3, 1.14 et 4.48~: 2.91, 1.17, 2.39 et \qty{3.67}{pW.kg^{-1}}.
\textbf{5.} \qty{10.1}{pW.kg^{-1}}, 2.07 fois celle d'aujourd'hui.
\textbf{6.} Potassium-40. Sa demi-vie, \qty{1.25}{Ga}, est courte~: plus de deux demi-vies
se sont écoulées depuis, et plus des trois quarts du \ce{^{40}K} de cette époque se sont désintégrés.
\textbf{7.} $1/(1 + 0.4\times2.7/4.567) = 0.809$~; $0.809\times1361 = \qty{1101}{W.m^{-2}}$.
\textbf{8.} $288\times0.809^{1/4} = \qty{273}{K}$, en dessous de \qty{283}{K}~: pas un état chaud valide~;
la Terre gèlerait.
\textbf{9.} $(1101\times0.4/(4\times0.616\,\sigma))^{1/4} = \qty{237}{K}$.
\textbf{10.} $0.616\times0.809 = 0.498$.
\textbf{11.} $0.004/\num{278e-6} = 14$ à $0.03/\num{278e-6} = 108$ fois.
\textbf{12.} Le thermostat de l'altération des silicates~: l'altération ralentit dans un climat froid, laissant
le dioxyde de carbone volcanique s'accumuler, et s'accélère dans un climat chaud.
\textbf{13.} 0.826 de l'actuel, en dessous de 0.932~: pas d'état chaud~; la Terre avec l'air actuel serait une
boule de neige.
\textbf{14.} $\frac{0.1}{0.9}\times1.519\times\num{1.962e7} = \qty{3.31}{Ma}$.
\textbf{15.} $\frac{0.2}{0.8}\times1.519\times\num{1.962e7} = \qty{7.45}{Ma}$.
\textbf{16.} Le \omterm{def:b2:precambrian-earth:snowball}{carbonate de couverture}, une dolomie laminée qui recouvre brutalement les dépôts glaciaires,
précipité dans une mer alimentée par une altération intense dans la phase chaude qui suit, et inondé par
l'eau de fonte.
\textbf{17.} Le fer(II) ne reste dissous que dans une eau sans oxygène~: le fer s'accumule dans
l'océan anoxique sous la glace et précipite en fer(III) là où l'oxygène arrive, par exemple avec
l'eau de fonte.
\textbf{18.} $(1616 - 1450)/(32.2 - 23) = \qty{18.0}{\celsius}$ par point de pourcentage.
\textbf{19.} $1616 + 18.0\times(30 - 32.2) = \qty{1576}{\celsius}$.
\textbf{20.} \qty{1616}{\celsius}, le point mesuré.
\textbf{21.} $1616 - 18.0\times20.2 = \qty{1252}{\celsius}$.
\textbf{22.} $1616 - 20\times20.2 = \qty{1212}{\celsius}$~: \qty{40}{\celsius} de moins, un dixième de
l'écart entre les laves~; la conclusion ne dépend pas du choix.
\textbf{23.} Les \omterm{def:b2:precambrian-earth:greenstone}{komatiites} venaient de panaches, plus chauds que le manteau ambiant, et de forts \omterm{def:b2:mantle-melting:fraction}{taux de
fusion} en profondeur~; la température d'une lave n'est pas la température potentielle de sa source. Le
manteau dans son ensemble était plus chaud de 200 à \qty{300}{K}.
\textbf{24.} La lave de \omterm{def:b2:precambrian-earth:greenstone}{komatiite} la plus chaude était d'environ \textbf{\qty{364}{\celsius}} plus chaude qu'un
basalte primaire des dorsales actuelles.
\end{solution}
